\section{Related Work}

% \textbf{TODO}

% Related Work (It could be cool to merge this into following sections instead of just having it at the beginning)

% SWA, Linear Attention, BigBird, etc. [accompany with experiments] -> These all make stuff faster, but they trade-off expressivity for speed, and continue to depend on full attention for performance

% Shallow Benchmarks: Most work focuses on eval on stuff that doesn’t require a lot of expressivity (HELMET). Point to a detailed section of the appendix with FAQs on benchmarks (e.g. Browsecomp, BRIGHT, etc.) we feel fit into this

% Better Benchmarks: Call out newer benchmarks (OOLONG, LongBench, etc.), discuss how they can be related, and how our analysis brings in a more careful selection of these tasks, while providing an explanation for the key thing that makes these more / less difficult. Also call out the fact that these benchmarks typically focus on pure evaluation (without in-domain training data), while we have scalable and controllable in-domain train/test settings. These also often don’t emphasize crucial differences in scaling behavior

% \prasannnote{RELATED WORK IS OLD}

% \paragraph{Long Context Stuff} \citep{Acharya2024StarAE, Beltagy2020LongformerTL}, motivated by prohibitive costs in LCLM training and inference, but relatively less work investigating how to maximally leverage expressivity of LCLMs. Architectural differences between different approaches, in particular chunked vs full attention, has been heavily explored in prior work. \citep{Ivgi2022SLED, Lu2024TurboRAGAR, Xiao2025EfficientLCICL} have all found that, with the right configuration, chunked attention seems to be able to perform similarly to full attention on various benchmarks. Crucially, however, the benchmarks used in this work, even ``multi-hop'' benchmarks which intuitively seem like they may require multi-step reasoning, generally don't require intensive inter-document reasoning (this mirrors the results we ourselves find on simpler tasks), and so doesn't explore the more sophisticated cases we examine in our work. This follows a general trend, with popular benchmarks such as HELMET \citep{Yen2024HELMETHT}, NarrativeQA \citep{Kocisk2017TheNR}, etc. all generally not requiring large-scale inter-document reasoning, potentially explaining shifts towards greater sparsity. \citep{Lee2024CanLL} mentions comparing placing queries before context in their paper, mentioning it to perform worse, however it's unclear if this is with format-specific training or not, which we find makes a big difference.

% \paragraph{Retrieval} Information retrieval is a long-standing field in computer science. Yet while progress has been made with dense approaches like DPR, and ColBERT, allowing for efficient search over large corpora, understanding of the limits of retrieval are somewhat less understood, in large part as evaluation is often confounded by high word overlap between query and relevant document, misalignment between train / test domain, and high correlations with existing world knowledge (e.g. BRIGHT's \cite{Su2024BRIGHTAR} CQA baseline is extremely high even without any retrieval). While these approaches will likely remain efficient and scalable solutions for simpler tasks, we hypothesize that for many more complex inter-document reasoning tasks, even approaches like deep research and agentic RAG, as a function of being built on these systems, will be fundamentally limited. While work exists hinting at these potential limitations \cite{Weller2025OnTT}, we believe our work provides additional and more concrete evidence towards long-term limitations of simple retrieval systems.

% \prasannnote{detailed pass up to here, except for abstract}

\paragraph{Prior Corpus Reasoning Tasks} Corpus reasoning tasks (including examples with possibly high-\cname) have been studied for decades in computer science, from various low-\cname tasks in our suite (e.g. IR \cite{TREC2021}), to others like event re-ordering (REVIEWER CITE). Fields like Exploratory Data Analysis \citep{Tukey1962TheFO} (which studies visualization and  understanding of large data), though often focused on structured data, have explored engineering specialized pipelines \cite{Shankar2024DocETLAQ} to understand unstructured data. Work in library sciences has also developed theories of interrelatedness between documents, most famously in the taxonomy of \citet{Tillett1991ATO}. While many of the settings we construct are new, we are not the first to propose tasks classifiable as high-\cname. That said, our \cname definition enables us to distinguish and understand what makes corpus reasoning tasks challenging. Consequently, to our knowledge this work is the first to identify LCLMs as uniquely suitable for high \cname tasks, and to study LCLMs for these tasks on growing corpus sizes---allowing us to identify limitations of the current low-\cname dominated regime. Our code, train data, and evaluation also provide the first large-scale (can generate contexts up to 10M+ tokens) testbed to study high-\cname tasks. 

\paragraph{Scalable Corpus Methods} Work applying language models to corpora has taken several directions. One path is to develop agentic scaffolds that combine efficient tools like dense retrievers \cite{Khattab2020ColBERTEA, Karpukhin2020DensePR} with LLMs. These are often hampered by fundamental challenges in retrieval \citep{Weller2025OnTT, Su2024BRIGHTAR, Wei2025BrowseCompAS}, but as shown by \cite{Zhang2025RecursiveLM} can be promising given the right scaffold and sufficient inference compute. Note that  the slowness of auto-regression (and the high per-query cost) make agentic systems intractable for high-complexity tasks---finding subtle contradictions via enumeration across $N$ documents may require generating on the order of $N^2$ tokens per query. Consequently, a complementary line of work---which we focus on in this work---seeks to develop architectures capable of ingesting full corpora end-to-end, often by making attention more efficient \citep{Acharya2024StarAE, Beltagy2020LongformerTL, Ivgi2022SLED, Lu2024TurboRAGAR, Xiao2025EfficientLCICL, Izacard2020LeveragingPR, Zaheer2020BigBT, Katharopoulos2020TransformersAR, Yang2024GatedDN}. Context extension (YarN), and length generalization (random rope) also fall into these efforts. Our work helps develop a unified understanding of what these can / can't achieve as a function of their computation. Importantly, our findings question some of this investigation, and reveal limitations obscured by the fact that these methods have primarily been developed on low-\cname tasks.

% \textbf{Information Retrieval} Information retrieval (IR) \cite{TREC2021} is a long-standing field, with many advances \cite{Khattab2020ColBERTEA}, and ongoing efforts to overcome fundamental limitations \cite{Weller2025OnTT} and progress on harder tasks \cite{Su2024BRIGHTAR, Wei2025BrowseCompAS}. While we classify many IR tasks as \ctc{N} \cname, and related methods even with ``agentic scaffolds'' \citep{Zhang2025RecursiveLM}, may be less suited to \ctc{N^2} tasks (hence why this work focuses on LCLMs), IR literature has developed many effective techniques \citep{HarPeled2012ApproximateNN} for scaling computation over corpora, and that may enable inform future work scaling \ctc{N^2} tasks.

% \textbf{Scaling LCLMs}  LCLMs are expensive, and much work  has investigated making them cheaper with attention mask or other architectural modifications while preserving popular LCLM benchmark performance (our attention mask approaches are based on this prior work) \citep{Yen2024HELMETHT, Yang2018HotpotQAAD, Kocisk2017TheNR}, but with one caveat---most popular tasks are \ctc{N}. Cheap dense retrieval systems \citep{Karpukhin2020DensePR, Khattab2020ColBERTEA} or scaffolds built on top of them \citep{Zhang2025RecursiveLM} \emph{are already well-suited to \ctc{N} tasks}, thus while \cmc{N} LCLM approaches may suffice on \ctc{N} tasks (note at large enough scale even this is non-trivial \cite{Lee2024CanLL}), we argue that LCLMs are best motivated for understudied \ctc{N^2} tasks, and encourage more work here. In a separate vein, linear attention \cite{Katharopoulos2020TransformersAR} and GDN \cite{Yang2024GatedDN} approaches have been motivated with the aim of much greater scalability, though still sometimes struggling on \ctc{N} NIAH settings \cite{Hsieh2024RULERWT}. \cite{Goldman2024IsIR} discusses how tasks beyond retrieval may involve different sorts of computation. \citep{Zhang2025RecursiveLM}, which argues for using agentic systems over LCLMs, informally mentions a notion of task complexity, including a pair enumeration task (which decomposes to a \ctc{N} pool creation steps followed by exhaustive enumeration), though doesn't relate this to architectural decisions or training, and focuses on less natural settings where REPL-based environment and command execution are suitable for many operations.  %\prasannnote{Hmm fix RLM reference to sound less defensive}% \prasannnote{Can't tell if this is better here or in related work. I think it's an important point so I feel like it should show up somewhere that's not at the end}



% \textbf{Older Work} Exploratory data analysis (EDA) is a long-standing field \citep{Tukey1962TheFO} investigating ways to better visualize and understand large quantities of data, but even modern approaches can rely on sophisticated / specialized pipelines \cite{Shankar2024DocETLAQ}, more structured (e.g. tabular) data, and typically operate below \ctc{N^2} \cname. Work in library sciences has also developed theories of interrelatedness between documents, most famously in the taxonomy of \citet{Tillett1991ATO}; our tasks with greater than \ctc{N} require leveraging one or more of these bibliographical relationships between documents. \prasann{TODO event re-ordering thing. Note I checked and there's barely any data so the thing reviewer sent isn't usable for us...}

